
Donc @@M0@@, qui fait de @@M1@@ un @@M2@@-torseur sur @@M3@@, est affine. D'ailleurs on sait (SGA 3 1.4.2) que @@M4@@ est une immersion fermée, donc affine, donc il en est de même du composé @@M5@@.

\textbf{Corollaire 1.5.} \uline{Soit} @@M6@@ \uline{un corps.}
\begin{enumerate}[label=\alph*)]
\item \uline{Soit} @@M7@@ \uline{un} @@M8@@\uline{-schéma propre. Alors} @@M9@@ \uline{est représentable} (J. P. Murre).
\item \uline{Soit} @@M10@@ \uline{un morphisme surjectif entre deux} @@M11@@\uline{-schémas propres. Alors le morphisme induit} @@M12@@ \uline{est affine.}
\end{enumerate}

\textbf{Corollaire 1.6.} \uline{Soient} @@M13@@ \uline{un schéma noethérien intègre,} @@M14@@ \uline{un} @@M15@@\uline{-morphisme surjectif entre deux} @@M16@@\uline{-schémas propres.}
\begin{enumerate}[label=\alph*)]
\item \uline{Si} @@M17@@ \uline{est normal sur} @@M18@@ \uline{(c'est-à-dire plat à fibres géométriquement normales), alors, quitte à restreindre} @@M19@@, @@M20@@ \uline{est quasi-fini.}
\item \uline{Si} @@M21@@ \uline{est essentiellement propre sur} @@M22@@ \uline{(c'est-à-dire, s'il satisfait au critère valuatif de propreté, par exemple, si} @@M23@@ \uline{est lisse sur} @@M24@@\uline{) alors quitte à restreindre} @@M25@@, @@M26@@ \uline{est fini (et en particulier, affine).}
\end{enumerate}

En effet dans a), pour tout @@M27@@, @@M28@@ est un groupe algébrique affine~; sa composante neutre est un groupe algébrique propre sur @@M29@@ (FGA 236, Théorème 2.1 (iii))~; @@M30@@ est donc fini. Dans b) @@M31@@ est à la fois propre et quasi-affine, donc fini.

\textbf{Remarque 1.7.} On ignore si 1.1 reste valable quand on y remplace « quasi-affine » par affine, comme pourraient le suggérer 1.5 et 1.6.

